\documentclass[11pt]{article}
\usepackage{progtable}

\begin{document}
\progheader{Sixième}{6e}

\begin{proglist}

\seq{6C01}{Fractions}{NUM}{Découvrir la notion de fraction et l'encadrement}{
\item 6C01-NUM01 — Connaître le vocabulaire des fractions (numérateur, dénominateur)
\item 6C01-NUM02 — Comparer une fraction avec l’unité
\item 6C01-NUM03 — Écrire une fraction sous forme décomposée
\item 6C01-NUM04 — Encadrer une fraction avec < et >
\item 6C01-NUM05 — Encadrer une fraction avec ≤ et ≥
\item 6C01-NUM06 — Placer une fraction sur une droite graduée
}

\seq{6C02}{Points et droites}{GEO}{Poser le vocabulaire de base en géométrie plane}{
\item 6C02-GEO01 — Reconnaître points et droites
\item 6C02-GEO02 — Utiliser le vocabulaire d’appartenance et d’alignement
\item 6C02-GEO03 — Reconnaître et tracer des demi-droites
\item 6C02-GEO04 — Construire un segment, sa longueur et son milieu
\item 6C02-GEO05 — Coder et lire une figure géométrique
\item 6C02-GEO06 — Distinguer droites sécantes et parallèles
\item 6C02-GEO07 — Reconnaître des droites perpendiculaires
\item 6C02-GEO08 — Tracer parallèles et perpendiculaires aux instruments
\item 6C02-GEO09 — Connaître les propriétés des droites parallèles
}

\seq{6C03}{Fractions égales}{NUM}{Reconnaître et produire des fractions égales}{
\item 6C03-NUM07 — Reconnaître des fractions égales
\item 6C03-NUM08 — Simplifier une fraction
}

\seq{6C04}{Angles}{GEO}{Utiliser le rapporteur avec précision}{
\item 6C04-GEO10 — Connaître le vocabulaire des angles
\item 6C04-GEO11 — Mesurer un angle au rapporteur
\item 6C04-GEO12 — Utiliser le rapporteur pour construire un angle
\item 6C04-GEO13 — Construire un angle de mesure donnée
\item 6C04-GEO14 — Construire la bissectrice d’un angle
}

\seq{6C05}{Nombres décimaux}{NUM}{Consolider la numération décimale}{
\item 6C05-NUM09 — Reconnaître une fraction décimale
\item 6C05-NUM10 — Lire et écrire un nombre décimal (écriture à virgule)
\item 6C05-NUM11 — Décomposer un nombre décimal (numération décimale)
\item 6C05-NUM12 — Placer un décimal sur une demi-droite graduée
\item 6C05-NUM13 — Comparer et encadrer des décimaux
\item 6C05-NUM14 — Donner une valeur approchée ou arrondie
}

\seq{6C06}{Cercles et triangles}{GEO}{Maîtriser le compas et le vocabulaire des triangles}{
\item 6C06-GEO15 — Tracer un cercle au compas
\item 6C06-GEO16 — Caractériser cercle et disque
\item 6C06-GEO17 — Construire des arcs et des figures avec des cercles
\item 6C06-GEO18 — Connaître le vocabulaire des polygones et des triangles
\item 6C06-GEO19 — Construire un triangle
\item 6C06-GEO20 — Connaître et utiliser l’inégalité triangulaire
\item 6C06-GEO21 — Reconnaître les triangles usuels
\item 6C06-GEO22 — Reconnaître les quadrilatères usuels
}

\seq{6C07}{Addition, soustraction, multiplication}{CALC}{Maîtriser le calcul décimal posé et les priorités opératoires}{
\item 6C07-CALC01 — Additionner et soustraire des nombres décimaux
\item 6C07-CALC02 — Multiplier un décimal par un entier
\item 6C07-CALC03 — Multiplier par 10, 100, 1000
\item 6C07-CALC04 — Multiplier deux nombres décimaux
\item 6C07-CALC05 — Appliquer les règles de priorité des opérations
}

\seq{6C08}{Gestion des données}{STAT}{Organiser, représenter et interpréter des données}{
\item 6C08-STAT01 — Lire et compléter un tableau de données
\item 6C08-STAT02 — Construire et lire un diagramme en barres
\item 6C08-STAT03 — Construire et lire un graphique
\item 6C08-STAT04 — Interpréter des données statistiques
}

\seq{6C09}{Positions d’angles}{GEO}{Reconnaître les positions relatives d’angles}{
\item 6C09-GEO23 — Reconnaître des angles adjacents
\item 6C09-GEO24 — Reconnaître des angles supplémentaires
\item 6C09-GEO25 — Reconnaître des angles opposés par le sommet
\item 6C09-GEO26 — Reconnaître des angles correspondants et alternes-internes
\item 6C09-GEO27 — Connaître la somme des angles d’un triangle
}

\seq{6C10}{Division entière}{CALC}{Maîtriser la division entière et la divisibilité}{
\item 6C10-CALC06 — Poser et effectuer une division entière
\item 6C10-CALC07 — Reconnaître multiples et diviseurs
\item 6C10-CALC08 — Connaître les critères de divisibilité par 2, 5, 10
\item 6C10-CALC09 — Connaître les critères de divisibilité par 3, 9
}

\seq{6C11}{Médiatrices}{GEO}{Construire médiatrices et cercle circonscrit}{
\item 6C11-GEO28 — Construire la médiatrice d’un segment
\item 6C11-GEO29 — Construire les médiatrices d’un triangle
\item 6C11-GEO30 — Construire le cercle circonscrit à un triangle
}

\seq{6C12}{Division décimale}{CALC}{Maîtriser la division décimale et la résolution de problèmes}{
\item 6C12-CALC10 — Poser et effectuer une division décimale
\item 6C12-CALC11 — Diviser par 10, 100, 1000
\item 6C12-CALC12 — Résoudre un problème en choisissant les bonnes opérations
}

\seq{6C13}{Aire et périmètre}{MEAS}{Distinguer et calculer périmètre et aire}{
\item 6C13-MEAS01 — Connaître les unités de longueur usuelles
\item 6C13-MEAS02 — Calculer le périmètre d’un polygone
\item 6C13-MEAS03 — Calculer le périmètre d’un disque
\item 6C13-MEAS04 — Connaître les unités d’aire usuelles
\item 6C13-MEAS05 — Calculer l’aire d’un rectangle ou d’un triangle rectangle
\item 6C13-MEAS06 — Calculer l’aire d’un triangle quelconque ou d’un disque
}

\seq{6C14}{Proportionnalité}{PROP}{Raisonner dans des situations proportionnelles}{
\item 6C14-PROP01 — Reconnaître des grandeurs proportionnelles
\item 6C14-PROP02 — Compléter un tableau de proportionnalité
\item 6C14-PROP03 — Utiliser les opérations sur les colonnes (linéarité)
\item 6C14-PROP04 — Utiliser le coefficient de proportionnalité
\item 6C14-PROP05 — Utiliser une échelle
\item 6C14-PROP06 — Tracer un diagramme circulaire
}

\seq{6C15}{Fractions et quotient}{NUM}{Relier fraction et quotient, additionner des fractions}{
\item 6C15-NUM15 — Reconnaître l’écriture fractionnaire d’un quotient
\item 6C15-NUM16 — Comparer et additionner des fractions de même dénominateur
\item 6C15-NUM17 — Comparer des fractions de dénominateurs quelconques
\item 6C15-NUM18 — Additionner des fractions de dénominateurs quelconques
}

\seq{6C16}{Symétrie axiale}{GEO}{Construire et caractériser la symétrie axiale}{
\item 6C16-GEO31 — Connaître la définition de la symétrie axiale
\item 6C16-GEO32 — Construire le symétrique d’une figure
\item 6C16-GEO33 — Connaître les propriétés de la symétrie axiale
\item 6C16-GEO34 — Construire le symétrique des figures usuelles
\item 6C16-GEO35 — Reconnaître les axes de symétrie d’une figure
}

\seq{6C17}{Proportions et pourcentages}{PROP}{Appliquer et calculer proportions et pourcentages}{
\item 6C17-PROP07 — Appliquer une proportion
\item 6C17-PROP08 — Calculer une proportion
\item 6C17-PROP09 — Appliquer un pourcentage
\item 6C17-PROP10 — Calculer un pourcentage
}

\seq{6C18}{Géométrie dans l’espace}{GEO}{Se repérer dans l’espace et calculer des volumes}{
\item 6C18-GEO36 — Reconnaître un assemblage de cubes en perspective
\item 6C18-GEO37 — Connaître le vocabulaire des solides usuels (cube, pavé, prisme, pyramide, cylindre, cône, boule)
\item 6C18-GEO38 — Connaître les unités de volume usuelles (cm³)
\item 6C18-GEO39 — Calculer le volume d’un cube ou d’un pavé droit
\item 6C18-GEO40 — Relier volume et contenance
}

\seq{6C19}{Probabilités}{STAT}{Découvrir le vocabulaire et le calcul des probabilités}{
\item 6C19-STAT05 — Décrire les issues d’une expérience aléatoire
\item 6C19-STAT06 — Calculer la probabilité d’un événement
\item 6C19-STAT07 — Distinguer des événements contraires
\item 6C19-STAT08 — Étudier la répétition d’expériences aléatoires
}

\seq{6C99}{Numération entière (non traité cette année)}{NUM}{Maîtriser la numération de position}{
\item 6C99-NUM19 — Lire et écrire les grands nombres entiers
\item 6C99-NUM20 — Décomposer un entier (unités, dizaines…)
\item 6C99-NUM21 — Comparer et ordonner des entiers
}

\end{proglist}
\end{document}
